\documentclass[12pt, a4paper]{book}

\usepackage[T1]{fontenc}
\usepackage[utf8]{inputenc}
\usepackage[charter, cal=cmcal]{mathdesign}
\usepackage[scaled=0.92]{helvet}
\usepackage{geometry}
\geometry{a4paper, top=30mm, bottom=30mm, inner=32mm, outer=25mm, marginparwidth=18mm}

\usepackage{amsmath, amsthm}
\usepackage[version=4]{mhchem}
\usepackage{graphicx}
\usepackage{xcolor}
\usepackage{tikz}
\usetikzlibrary{arrows.meta, shapes.geometric, positioning, calc, backgrounds}
\usepackage{pgfplots}
\pgfplotsset{compat=1.18}
\usepackage{tcolorbox}
\tcbuselibrary{skins, breakable, theorems}
\usepackage{titlesec}
\usepackage{fancyhdr}
\usepackage{enumitem}
\usepackage{booktabs}
\usepackage{array}
\usepackage{parskip}
\usepackage{siunitx}
\usepackage[colorlinks=true, linkcolor=BrickRed,
            urlcolor=NavyBlue, citecolor=ForestGreen]{hyperref}

% ── Colours (Charcoal / Royal Blue) ──────────────────────────────────────────
\definecolor{JourneyColor}{HTML}{0F172A}
\definecolor{ModuleColor}{HTML}{2563EB}
\definecolor{PuzzleColor}{HTML}{D97706}
\definecolor{DiscoveryColor}{HTML}{3B82F6}
\definecolor{WarningColor}{HTML}{991B1B}
\definecolor{ProblemColor}{HTML}{6B21A8}
\definecolor{ReflectColor}{HTML}{0D9488}
\definecolor{HistoryColor}{HTML}{78350F}

% ── Boxes ─────────────────────────────────────────────────────────────────────
\newtcolorbox{puzzlebox}[1][]{enhanced,breakable,
  colback=PuzzleColor!6,colframe=PuzzleColor,fonttitle=\bfseries\sffamily,
  title={The Puzzle},attach boxed title to top left={yshift=-2mm,xshift=6mm},
  boxed title style={colback=PuzzleColor},borderline west={3pt}{0pt}{PuzzleColor},#1}

\newtcolorbox{discoverybox}[1][]{enhanced,breakable,
  colback=DiscoveryColor!8,colframe=DiscoveryColor,fonttitle=\bfseries\sffamily,
  title={Discovery},attach boxed title to top left={yshift=-2mm,xshift=6mm},
  boxed title style={colback=DiscoveryColor},borderline west={3pt}{0pt}{DiscoveryColor},#1}

\newtcolorbox{problembox}[1][]{enhanced,breakable,
  colback=ProblemColor!5,colframe=ProblemColor,fonttitle=\bfseries\sffamily,
  title={Problems},attach boxed title to top left={yshift=-2mm,xshift=6mm},
  boxed title style={colback=ProblemColor},borderline west={3pt}{0pt}{ProblemColor},#1}

\newtcolorbox{deepbox}[1][]{enhanced,breakable,
  colback=JourneyColor!5,colframe=JourneyColor,fonttitle=\bfseries\sffamily,
  title={Deep Problems},attach boxed title to top left={yshift=-2mm,xshift=6mm},
  boxed title style={colback=JourneyColor},borderline west={4pt}{0pt}{JourneyColor},#1}

\newtcolorbox{historybox}[1][]{enhanced,breakable,
  colback=HistoryColor!7,colframe=HistoryColor!60,fonttitle=\bfseries\sffamily\small,
  title={History Note},attach boxed title to top left={yshift=-2mm,xshift=6mm},
  boxed title style={colback=HistoryColor!80},#1}

\newtcolorbox{reflectionbox}[1][]{enhanced,breakable,
  colback=ReflectColor!6,colframe=ReflectColor,fonttitle=\bfseries\sffamily,
  title={Reflection},attach boxed title to top left={yshift=-2mm,xshift=6mm},
  boxed title style={colback=ReflectColor},borderline west={3pt}{0pt}{ReflectColor},#1}

\newtcolorbox{namedconcept}[2][]{enhanced,
  colback=ModuleColor!6,colframe=ModuleColor,fonttitle=\bfseries\sffamily\large,
  title={#2},attach boxed title to top left={yshift=-2mm,xshift=6mm},
  boxed title style={colback=ModuleColor},#1}

% ── Section formatting ────────────────────────────────────────────────────────
\titleformat{\chapter}[display]
  {\normalfont\sffamily\huge\bfseries\color{JourneyColor}}
  {\textcolor{JourneyColor!60}{\chaptertitlename\ \thechapter}}{10pt}{\Huge}
  [\vspace{2pt}\textcolor{JourneyColor}{\titlerule[1.5pt]}]
\titleformat{\section}
  {\normalfont\sffamily\Large\bfseries\color{ModuleColor}}
  {\thesection}{1em}{}[\textcolor{ModuleColor!40}{\titlerule[0.6pt]}]
\titleformat{\subsection}
  {\normalfont\sffamily\large\bfseries\color{JourneyColor!80}}{\thesubsection}{1em}{}
\titleformat{\subsubsection}
  {\normalfont\sffamily\normalsize\bfseries}{\thesubsubsection}{1em}{}

\pagestyle{fancy}
\fancyhf{}
\fancyhead[LE]{\small\sffamily\textcolor{JourneyColor}{\leftmark}}
\fancyhead[RO]{\small\sffamily\textcolor{ModuleColor}{\rightmark}}
\fancyfoot[C]{\small\thepage}
\renewcommand{\headrulewidth}{0.4pt}

\newcommand{\newterm}[1]{\textit{\textcolor{DiscoveryColor}{#1}}}
\newcommand{\important}[1]{\textcolor{ModuleColor}{\textbf{#1}}}

\begin{document}

\setcounter{chapter}{5}
\chapter[Journey 6 --- Distributed Systems and Concurrency]
        {Journey Six\\[8pt]
         \textcolor{ModuleColor}{\large Distributed Systems and Concurrency}}
\setcounter{section}{2}

\section{Module 6.3 --- Distributed Systems, Consensus, and Paxos}
\label{sec:mod6.3}

\subsection*{Motivation}

We completed in Module 6.2 the physical memory boundaries of concurrency, studying store buffers, cache coherence (MESI), and memory fences.
Now, we scale out past CPU bus connections to network boundaries.
Instead of cores coordinating over a shared memory bus, we must analyze independent servers communicating over unreliable network channels.
We enter the domain of **Distributed Systems**.

In this module, we formalize the fundamental limits of network computation.
We will study the **CAP Theorem**, proving why systems cannot simultaneously guarantee consistency and availability during network partitions.
We will analyze the consensus problem, state the **FLP Impossibility Theorem**, and explore the mechanics of **Paxos**, deriving how majority quorums enforce safe state consensus in the presence of node crash failures.

\subsection{The Unreliable Channel}

We investigate network partition trade-offs and consensus limitations.

\begin{puzzlebox}
\textbf{The Distributed Database.}

Suppose we have a distributed database running on three servers ($G_1, G_2, G_3$) located in different cities.
Clients can write data to any server, and the servers replicate updates by sending network messages to each other.
\begin{enumerate}[label=(\alph*), leftmargin=*]
  \item Suppose a network failure partitions the servers: $G_1$ and $G_2$ can communicate, but their link to $G_3$ is cut.
        * A client writes a new value `x = 5` to server $G_1$.
        * Another client reads variable `x` from server $G_3$ concurrently.
        Explain why the system must choose between **Consistency** (preventing $G_3$ from returning stale data) and **Availability** (allowing both writes and reads to succeed without blocking).
  \item This limit is formalized by the **CAP Theorem**.
        Define the three parameters: Consistency (Linearizability), Availability, and Partition Tolerance.
        Prove that when a partition ($P$) occurs, we must choose between Consistency ($C$) and Availability ($A$).
  \item To maintain consistency under node crashes, servers must agree on a single sequence of updates. This is the **Consensus Problem**.
        We require three properties of any consensus algorithm:
        * **Agreement:** No two correct nodes decide on different values.
        * **Validity:** The decided value must have been proposed by some node.
        * **Termination:** Every correct node eventually decides.
        State the **FLP Impossibility Theorem** (1985). Why does even a single node crash prevent guaranteed termination in asynchronous networks?
\end{enumerate}
\end{puzzlebox}

\subsection{The Paxos Consensus Protocol}

We study the roles, phases, and quorum constraints of Paxos.

\begin{namedconcept}{Paxos Protocol Architecture}
Paxos guarantees safety (Agreement and Validity) under crash failures, bypassing the FLP termination limit by allowing infinite loops only in rare, non-terminating configurations (resolved in practice by randomized timeouts or leaders).
Nodes play three roles:
* **Proposers:** Advocate for values proposed by clients.
* **Acceptors:** Store and vote on proposals, acting as the memory of the consensus.
* **Learners:** Detect when a proposal has been decided (accepted by a majority).
\end{namedconcept}

\begin{explobox}
Paxos operates in two distinct phases:
* **Phase 1a (Prepare):** A Proposer selects a unique proposal number $n$ and broadcasts `Prepare(n)` to a majority of Acceptors.
* **Phase 1b (Promise):** An Acceptor receives `Prepare(n)`. If $n$ is greater than any proposal number it has seen, it promises not to accept any proposals numbered $< n$, and replies with the highest-numbered proposal it has already accepted (if any).
* **Phase 2a (Accept):** If the Proposer receives promises from a majority, it sends `Accept(n, v)` to those Acceptors, where $v$ is the value of the highest-numbered proposal among the promises, or its own proposed value if no accepted value was returned.
* **Phase 2b (Accepted):** An Acceptor receives `Accept(n, v)`. It accepts the proposal unless it has already promised not to.
\begin{enumerate}[label=(\alph*), leftmargin=*]
  \item Why must the proposal numbers $n$ be unique across all proposers? (What would happen if two proposers used the same $n$ for different values?)
  \item Explain how the **Majority Quorum** ($Q > N/2$) ensures that any two quorums must intersect.
        Why is this intersection property critical for carrying over previously accepted values to Phase 2?
  \item Draw a message sequence chart showing a successful Paxos round where a value $v$ is decided among three acceptors.
\end{enumerate}
\end{explobox}

\subsection{Fault Tolerance and Raft Consensus}

We analyze leader election and log replication.

\begin{explobox}
While Paxos is mathematically elegant, implementing it for a continuous stream of commands (a replicated log) is complex.
Modern systems often use **Raft**, which decomposes consensus into explicit leader election, log replication, and safety.
\begin{enumerate}[label=(\alph*), leftmargin=*]
  \item In Raft, time is divided into arbitrary **Terms** of numbered integers.
        During an election, a candidate must obtain votes from a majority of nodes to become leader.
        Prove that there can be at most one leader elected in any single Term.
  \item Raft enforces the **Leader Completeness** safety invariant: if a log entry is committed in a given term, that entry will be present in the logs of the leaders for all higher-termed terms.
        How does the voting process (where a follower denies its vote if the candidate's log is less up-to-date than its own) guarantee this completeness invariant?
  \item If a distributed database has $N = 5$ nodes, what is the maximum number of simultaneous node crash failures the system can tolerate while remaining functional?
\end{enumerate}
\end{explobox}

\newpage

% ─── Section: Discovery ──────────────────────────────────────────────────────
\section{Part 3 --- Discovery}

\begin{discoverybox}
\textbf{What have we built?}

\textbf{First.} The CAP Theorem proves that no distributed system can guarantee both Consistency and Availability during network partitions.

\textbf{Second.} The FLP Impossibility Theorem states that deterministic consensus cannot guarantee termination in asynchronous networks if a single node can crash.

\textbf{Third.} Paxos achieves safe consensus by using a two-phase prepare-accept protocol where majority quorums intersect, preserving chosen states.

\textbf{Fourth.} Replicated consensus engines (like Raft) structure Paxos concepts into terms, leader election, and log replication.
\end{discoverybox}

\newpage

% ─── Section: Naming Concepts ────────────────────────────────────────────────
\section{Part 4 --- What This Is Called}

\begin{namedconcept}{Distributed parameters}
* **CAP Theorem:** The theorem stating a distributed system can satisfy at most two of Consistency, Availability, and Partition Tolerance.
* **Consensus:** The process of reaching agreement among a set of independent network nodes on a shared data value.
* **FLP Impossibility:** The theorem proving asynchronous consensus cannot guarantee termination under crash failures.
* **Paxos:** A family of protocols for solving consensus in a network of unreliable or fallible processors.
* **Quorum:** The minimum number of node votes required to validate a transaction or decision.
* **Raft:** A leader-based consensus algorithm designed for ease of understanding and implementation.
\end{namedconcept}

\newpage

% ─── Section: Solved Exercises ───────────────────────────────────────────────
\section{Part 5 --- Solved Exercises}

\begin{problembox}
\begin{enumerate}[label=\textbf{E\arabic*.}]
  \item \textbf{Quorum Intersection Proof.}
  
        Prove that in a system of $N$ nodes, any two majority quorums must share at least one node.
        
        \textbf{Solution:}
        Let $V$ be the set of nodes, with size $|V| = N$.
        * A majority quorum is any subset $Q \subseteq V$ such that:
          \[
            |Q| \;\ge\; \left\lfloor \frac{N}{2} \right\rfloor + 1.
          \]
        Let $Q_1$ and $Q_2$ be two majority quorums.
        We want to show that $Q_1 \cap Q_2 \neq \emptyset$.
        
        1. **Sum the sizes of the quorums:**
           \[
             |Q_1| + |Q_2| \;\ge\; 2 \left( \left\lfloor \frac{N}{2} \right\rfloor + 1 \right).
           \]
           * If $N$ is even ($N = 2k$):
             \[
               |Q_1| + |Q_2| \;\ge\; 2(k + 1) \;=\; 2k + 2 \;=\; N + 2.
             \]
           * If $N$ is odd ($N = 2k+1$):
             \[
               |Q_1| + |Q_2| \;\ge\; 2(k + 1) \;=\; 2k + 2 \;=\; N + 1.
             \]
           In both cases, we have:
           \[
             |Q_1| + |Q_2| \;\ge\; N + 1 \;>\; |V|.
           \]
        2. **Apply the Principle of Inclusion-Exclusion:**
           The size of the union of the two sets is:
           \[
             |Q_1 \cup Q_2| \;=\; |Q_1| + |Q_2| - |Q_1 \cap Q_2|.
           \]
           Since $Q_1 \cup Q_2 \subseteq V$, we must have $|Q_1 \cup Q_2| \le |V| = N$.
           Substituting this in:
           \[
             N \;\ge\; |Q_1| + |Q_2| - |Q_1 \cap Q_2| \;\ge\; N + 1 - |Q_1 \cap Q_2|.
           \]
        3. **Isolate the intersection term:**
           \[
             |Q_1 \cap Q_2| \;\ge\; 1.
           \]
        Since the size of the intersection is at least 1, the two majority quorums must share at least one node.
        This completes the proof.
        
  \item \textbf{CAP Theorem Partition Trade-off Trace.}
  
        Sketch a formal counterexample showing why Consistency and Availability cannot coexist in a partitioned 2-node system.
        
        \textbf{Solution:}
        Let the system consist of two nodes, $G_1$ and $G_2$, holding a replicated variable $x$ (initially $x = 0$).
        A network partition cuts the link between $G_1$ and $G_2$.
        
        1. **Event 1: Write operation.**
           A client connects to $G_1$ and requests a write: `write(x = 5)`.
           * To be **Available**, $G_1$ must accept the write immediately and return a success acknowledgment to the client without waiting indefinitely.
           * $G_1$ updates its local value to $x = 5$.
           * Because of the partition, $G_1$ cannot send this update to $G_2$.
           
        2. **Event 2: Read operation.**
           Concurrently, another client connects to $G_2$ and requests a read: `read(x)`.
           * To be **Available**, $G_2$ must respond to the client immediately.
           * Since it has not received the update, $G_2$ returns its local value: $x = 0$.
           
        3. **Analyze the outcome:**
           * The read operation at $G_2$ returned $0$, which is stale compared to the completed write of $5$ at $G_1$. This violates **Consistency** (specifically, linearizability).
           * To enforce consistency, $G_2$ would have to block or return an error, which violates **Availability**.
        
        Thus, under a network partition, the system must choose to sacrifice either Consistency or Availability.
\end{enumerate}
\end{problembox}

\newpage

% ─── Section: Problems ───────────────────────────────────────────────────────
\section{Part 6 --- Problems}

\begin{problembox}
\begin{enumerate}[label=\textbf{P\arabic*.}]
  \item \textbf{Paxos Split Vote.}
  
        Suppose two Proposers propose values $v_1$ (with proposal number $n_1 = 1$) and $v_2$ (with proposal number $n_2 = 2$) concurrently to three Acceptors $\{A_1, A_2, A_3\}$ in a Paxos consensus network.
        Suppose the network messages interleave as follows:
        1. Proposer 1 sends `Prepare(1)` to $A_1$ and $A_2$. Both send promises.
        2. Proposer 2 sends `Prepare(2)` to $A_2$ and $A_3$. Both send promises.
        \begin{enumerate}[label=(\alph*)]
          \item When $A_2$ receives `Prepare(2)`, explain why it is allowed to send a promise to Proposer 2, even though it has already promised Proposer 1. (What happens to its old promise?)
          \item Proposer 1 now sends `Accept(1, v_1)` to $A_1$ and $A_2$.
                Explain why $A_1$ accepts the proposal, but $A_2$ rejects it.
          \item Trace the remaining steps of the protocol. Show which value ($v_1$ or $v_2$) is ultimately decided by a majority quorum.
        \end{enumerate}

  \item \textbf{Raft Split Elections.}
  
        In Raft, a split election occurs when multiple nodes candidate concurrently, and none obtains a majority of votes.
        \begin{enumerate}[label=(\alph*)]
          \item Consider a system of $N = 4$ nodes.
                Suppose Candidate 1 gets 2 votes, and Candidate 2 gets 2 votes.
                Explain why neither candidate can become leader.
          \item How does Raft resolve this split election block?
                *(Explain the use of randomized election timeouts).*
          \item If $N = 4$ nodes, why is the split vote probability higher than in a system of $N = 5$ nodes?
        \end{enumerate}
\end{enumerate}
\end{problembox}

\newpage

% ─── Section: Deep Problems ──────────────────────────────────────────────────
\section{Part 7 --- Deep Problems}

\begin{deepbox}
\begin{enumerate}[label=\textbf{D\arabic*.}]
  \item \textbf{Proof of Paxos Safety Invariant.}
  
        The core safety invariant of Paxos is: if a proposal with value $v$ and number $n$ is chosen (accepted by a majority of acceptors), then every proposal with number $n' > n$ issued by any proposer also has value $v$.
        
        Let $S$ be the majority quorum of Acceptors that accepted proposal $(n, v)$.
        Let $T$ be the majority quorum of Acceptors that respond to a Proposer's `Prepare(n')` message for a later proposal $n' > n$.
        
        \begin{enumerate}[label=(\alph*)]
          \item Using the quorum intersection theorem, prove that there exists at least one Acceptor $u$ such that $u \in S \cap T$.
          \item Since $u \in S$, it has accepted proposal $(n, v)$.
                Since $u \in T$, it must respond to the proposer of $n'$ with its accepted value.
                Explain how the proposer of $n'$ receives this value $v$ during the Phase 1b promise step.
          \item Under the Paxos Phase 2a rule, the proposer must set the value of its proposal $n'$ to the value of the highest-numbered proposal among all responses.
                Using mathematical induction on the proposal numbers, prove that the proposer of $n'$ is forced to choose value $v$, preserving the decided consensus.
        \end{enumerate}

  \item \textbf{FLP Impossibility and Asynchronous Time boundaries.}
  
        The FLP Impossibility Theorem states that in an asynchronous network model, no deterministic consensus protocol can be both safe (Agreement and Validity) and live (guaranteed Termination) under a single unannounced crash failure.
        
        \begin{enumerate}[label=(\alph*)]
          \item Contrast **Synchronous** networks (which have bounded message transmission delays) with **Asynchronous** networks (where message delays are unbounded but finite).
                Why does this distinction alter consensus solvability?
          \item In an asynchronous network, if a node has not responded to a message after 1 hour, explain why we cannot distinguish between:
                * The node has crashed.
                * The node is running extremely slowly or the network link is delayed.
                How does this ambiguity prevent deterministic protocols from achieving termination?
          \item Since Paxos is safe and widely used, how does it bypass the FLP impossibility result in practical deployments?
        \end{enumerate}
\end{enumerate}
\end{deepbox}

\newpage

% ─── Section: Reflection ─────────────────────────────────────────────────────
\section{Part 8 --- Reflection}

\begin{reflectionbox}
\textbf{Conceptual Reflection}

Consider the following questions:
\begin{enumerate}
  \item **On the bounds of consensus:** Paxos requires majorities to agree. This means a system can only tolerate $F < N/2$ crash failures. Is this a fundamental mathematical boundary of information, or can we design systems that tolerate more failures using alternative quorums?
  \item **On the nature of consistency:** Consistency (Linearizability) forces all nodes to behave as if there is only a single copy of the database. This restricts performance by requiring constant network communication. When is it acceptable to relax consistency (e.g. eventual consistency) to maximize database throughput?
  \item **On trust in networks:** Consensus protocols assume that crashed nodes stop executing (fail-stop model). In real networks, nodes can experience **Byzantine** failures: executing maliciously or sending corrupt messages. How does the consensus problem change when we must tolerate Byzantine nodes? (Think about blockchain systems).
\end{enumerate}
\end{reflectionbox}

\medskip
\noindent\textbf{Looking ahead.}
Module~6.4 introduces **Cryptography, Secure Computation, and Zero Knowledge**. Having analyzed how distributed nodes agree on shared data, we will study how to protect this data: symmetric/asymmetric encryption, digital signatures, zero-knowledge proofs (ZKPs), and secure multi-party computation.

\vspace{1.4cm}
\begin{center}
{\color{JourneyColor}\rule{0.5\linewidth}{0.8pt}}\\[6pt]
{\small\color{JourneyColor}\textit{End of Journey 6, Module 6.3 --- Distributed Systems, Consensus, and Paxos}}\\[4pt]
{\small\color{JourneyColor}Continue to: \textbf{Module 6.4 --- Cryptography, Secure Computation, and Zero Knowledge}}\\[6pt]
{\color{JourneyColor}\rule{0.5\linewidth}{0.8pt}}
\end{center}

\end{document}
